\section*{अध्याय \ref{ch:b3:cancer-biology} --- कर्कट जीवविज्ञान}

\begin{solution}{exo:b3:cancer-biology:1}
कोई \omterm{def:b3:cancer-biology:genes}{कर्कटजीन} ऐसे जीन का प्रकार्य-लाभ वाला रूप है जो प्रचुरोद्भवन या
उत्तरजीविता को बढ़ावा देता है; एक ही उत्परिवर्ती युग्मविकल्पी पर्याप्त
है (प्रभावी): \emph{RAS} (बिंदु उत्परिवर्तन), \emph{MYC} (स्थानांतरण,
प्रवर्धन), \emph{HER2}, \emph{BCR--ABL}। कोई अर्बुद-दमनकारी
प्रचुरोद्भवन रोकता है या विराम, मृत्यु अथवा मरम्मत थोपता है; दोनों
युग्मविकल्पी खोने पड़ते हैं (कोशिका-स्तर पर अप्रभावी): \emph{RB},
\emph{TP53}, \emph{APC}, \emph{BRCA1}।
\end{solution}

\begin{solution}{exo:b3:cancer-biology:2}
टिकाऊ प्रचुरोद्भवन: \emph{RAS}, \emph{MYC}। वृद्धि-दमनकारियों के प्रति
असंवेदिता: \emph{RB}, \emph{CDKN2A} (p16)। मृत्यु का प्रतिरोध:
\emph{BCL2}, \emph{TP53}। प्रतिकृतिक अमरता: \omterm{def:b3:cancer-biology:telomerase}{टीलोमरेज़} का
प्रवर्तक। रक्तवाहिनीजनन: HIF के ज़रिए प्रेरित \emph{VEGF} (\emph{VHL}
का ह्रास)। \omterm{def:b3:cancer-biology:metastasis}{घुसपैठ} और \omterm{def:b3:cancer-biology:hallmarks}{विक्षेपण}: E-कैडहेरिन का ह्रास,
उपकला--मध्यजनोतकी कारक। जीनोम अस्थिरता: असंगति-मरम्मत
जीन, \emph{\omterm{def:b3:dna-repair:dsb}{BRCA}}। अनियमित उपापचय: \emph{MYC}, \omterm{prop:b3:cancer-biology:pathways}{PI3K पथ}।
प्रतिरक्षा-अपवंचन: \emph{PD-L1}, MHC का ह्रास। अर्बुद-प्रवर्तक शोथ:
\emph{NF-$\kappa$B} संकेतन।
\end{solution}

\begin{solution}{exo:b3:cancer-biology:3}
जिन बच्चों का कुल-इतिहास था उनमें दोनों आँखों में कई \omterm{def:b3:cancer-biology:hallmarks}{अर्बुद} जल्दी बने;
छिटपुट स्थितियों में एक ही \omterm{def:b3:cancer-biology:hallmarks}{अर्बुद}, बाद में। उन्होंने निष्कर्ष निकाला कि
किसी कोशिका में दो घटनाएँ चाहिए: कुलगत रोगियों ने एक वंशागत पाई थी और
उन्हें केवल दूसरी, कायिक घटना चाहिए थी (जो दृष्टिपटल की कई कोशिकाओं में
होने जितनी बारंबार है), जबकि छिटपुट रोगियों को दोनों एक ही कोशिका में
कायिक रूप से चाहिए थीं --- दुर्लभ, देर से, अकेली। इसलिए ऐसा जीन जिसकी
दोनों प्रतियाँ खोनी पड़ें: पहला अर्बुद-दमनकारी।
\end{solution}

\begin{solution}{exo:b3:cancer-biology:4}
हिस्टिडीन-माँगता कोई \emph{Salmonella} उत्परिवर्ती परीक्षण-यौगिक के साथ
बिना हिस्टिडीन के फैलाया जाता है; बस्तियाँ केवल उन्हीं कोशिकाओं से उठती
हैं जिनमें किसी नए उत्परिवर्तन ने दोष पलट दिया हो, इसलिए उनकी संख्या
उत्परिवर्तजनकता मापती है। यकृत-निष्कर्ष इसलिए जोड़ा जाता है कि अनेक
कर्कटजन तब तक हानिरहित रहते हैं जब तक यकृत के एंज़ाइम उन्हें क्रियाशील
उपापचयजों में न बदल दें, जैसा शरीर में होता है; जीवाणुओं में वे एंज़ाइम
नहीं होते।
\end{solution}

\begin{solution}{exo:b3:cancer-biology:5}
आयु दोगुनी होते-होते घटना-दर $100$ गुना चढ़ती है: $2^{k-1} = 100$, $k - 1 =
\log_{2}100 = 6.6$, इसलिए भोले पाठ पर लगभग सात या आठ दर-सीमित घटनाएँ
--- आघातों के बीच क्लोनी फैलाव मानें तो कम।
\end{solution}

\begin{solution}{exo:b3:cancer-biology:6}
$L = \sqrt{2Dc_{0}/q}$: $q = 0.01$ के लिए $\sqrt{2\times 2\times 10^{-9}
\times 0.05/0.01} = \qty{141}{\micro m}$; $q = 0.03$ के लिए
\qty{82}{\micro m}। तेज़ी से बढ़ता \omterm{def:b3:cancer-biology:hallmarks}{अर्बुद} तेज़ी से श्वसन करता है, ऑक्सीजन
कम दूरी में चुका देता है, और उसका केंद्र छोटे आकार पर ही मर जाता है।
\end{solution}

\begin{solution}{exo:b3:cancer-biology:7}
$10^{11}\times 10^{-3n} < 1$ के लिए $n > 3.7$ चाहिए: चार बारियाँ। दो
बारियों के बाद $10^{5}$ कोशिकाएँ बचती हैं --- घन मिलीमीटर का दसवाँ
भाग, यानी उन $10^{9}$ से बहुत नीचे जो पकड़ में आ सकती हैं --- इसलिए हर
जाँच के हिसाब से \omterm{def:b3:cancer-biology:hallmarks}{अर्बुद} ``चला गया'' है जबकि एक लाख कोशिकाएँ बची हैं;
उपचार नियोजित बारियों तक चलना चाहिए, न कि \omterm{def:b3:cancer-biology:hallmarks}{अर्बुद} के लुप्त होने तक।
\end{solution}

\begin{solution}{exo:b3:cancer-biology:8}
दोनों घाव एक ही पथ सक्रिय करते हैं (ग्राही $\to$ \omterm{def:b3:cancer-biology:genes}{Ras} $\to$ ERK);
एक बार एक हो जाए तो दूसरा कोई और लाभ नहीं देता और चुना नहीं जाता।
कोई EGFR निरोधक ग्राही रोकता है; जो कोशिका नीचे की ओर कोई
KRAS उत्परिवर्तन पा ले वह ग्राही के रुके रहते ही संकेत बहाल कर देती है,
और उपचार के दौरान चुन ली जाती है --- बाईपास से प्रतिरोध।
\end{solution}

\begin{solution}{exo:b3:cancer-biology:9}
E7 Rb से E2F मुक्त कर देता है, जिससे संक्रमित उपकला कोशिका बिना वृद्धि
कारकों के \omterm{def:b3:cell-cycle-apoptosis:checkpoints}{निर्बंधन बिंदु} पार कर जाती है; E6 \omterm{def:b3:cancer-biology:genes}{p53} नष्ट कर देता है, इसलिए
अनुचित प्रचुरोद्भवन और उससे होने वाली क्षति न विराम चलाती है न
\omterm{def:b3:cell-cycle-apoptosis:apoptosis}{एपॉप्टोसिस}। दो अभिलक्षण एक साथ मिल जाते हैं, पर शेष --- अमरता, \omterm{def:b3:cancer-biology:metastasis}{घुसपैठ},
अस्थिरता --- के लिए और कायिक उत्परिवर्तन चाहिए, जिन्हें टिकाऊ रूप से
संक्रमित कोशिकाओं में जमा होने में दशक लगते हैं। टीका सम्पुट के
विरुद्ध उदासीनकारी प्रतिरक्षी उठाता है और संक्रमण रोक देता है; संपर्क
से पहले दिया जाए तो रूपांतरित होने को कोई संक्रमित कोशिका ही नहीं होती।
\end{solution}

\begin{solution}{exo:b3:cancer-biology:10}
$i$-वें आघात के बाद क्लोन $g$ से बढ़ता है, इसलिए अगले आघात के जोखिम
वाली कोशिकाएँ $g^{i}$ गुना हो जाती हैं; किसी एक मूल कोशिका से उतरे
वंशक्रम में पूरे अनुक्रम की प्रायिकता $g\cdot g^{2}\cdots$ है
--- प्रति आघात अचर $g$ के लिए संचयी जोखिम क्रम-निर्भर गुणकों तक
$N(ut)^{k}
g^{k-1}$ बन जाता है, और घटना-दर $t^{k-1}$ घात रखती है पर पूर्वगुणक
$g^{k-1}$ गुना बड़ा। यदि फैलाव स्वयं समय के साथ बढ़ें (चरघातांकी रूप से
बढ़ते क्लोन), तो वक्र और तीखा हो जाता है, इसलिए $6$ की ढाल सात से कम
चालकों से बन जाती है: आर्मिटेज--डॉल का $k$ दर-सीमित चालकों की संख्या की
ऊपरी सीमा है, और अनुक्रमित \omterm{def:b3:cancer-biology:hallmarks}{अर्बुद} दो से आठ दिखाते हैं।
\end{solution}

\begin{solution}{exo:b3:cancer-biology:11}
(क) अकेली औषध A पहले से मौजूद $\mu_{A}N = 100$ प्रतिरोधी कोशिकाओं का
सामना करती है: $P_{0} = e^{-100} \approx 0$; A-प्रतिरोधी क्लोन बच जाता
है और A के अठारह सप्ताहों में फिर बढ़ जाता है, और B के शुरू होने तक वह
फिर कोई बड़ी समष्टि है, इसलिए B भी $\mu_{B}N' \gg 1$ का सामना करती है
--- क्रम दो बार विफल होता है। (ख) साथ: कोई दोहरा-प्रतिरोधी कोशिका प्रति
विभाजन $10^{-14}$ से उठती है, $\mu_{A}\mu_{B}N = 10^{-5}$, $P_{0} =
0.99999$। आरंभ से संयोजन, जब $N$ सबसे छोटा हो, इसलिए नियम है कि दो
छोटी दरों का गुणनफल ही प्रतिरोध को असंभाव्य बनाता है।
\end{solution}

\begin{solution}{exo:b3:cancer-biology:12}
$100$ गुना अधिक कोशिकाओं और उन्हीं $u$ तथा $k$ के साथ संचयी जोखिम
$N(ut)^{k}$ $100$ गुना ऊँचा होता --- हर हाथी को जवानी में \omterm{def:b3:cancer-biology:hallmarks}{कर्कट} से मरना
चाहिए। ऐसा होता नहीं, इसलिए $u$ या $k$ भिन्न होने चाहिए:
\emph{TP53} की बीस प्रतियाँ क्षति पर एपॉप्टोटिक अनुक्रिया प्रबल कर
देती हैं और उत्परिवर्ती कोशिकाएँ फैलने से पहले हटा देती हैं (उस चरण के
लिए प्रभावी $u$ घटाकर); बड़े प्राणियों में प्रति कोशिका उत्परिवर्तन-दर
भी कम होती है (बेहतर मरम्मत, प्रति कोशिका प्रति वर्ष कम विभाजन,
क्योंकि उनकी कोशिकाएँ धीरे बदलती हैं), और उन्हें अधिक आघात चाहिए हो
सकते हैं (दमनकारियों की एक अतिरिक्त परत)। कर्कट-दमन शरीर-आकार के साथ
विकसित हुआ।
\end{solution}

\begin{solution}{pb:b3:cancer-biology:1}
\textbf{1.} $10^{12}/10^{3} = 10^{9}$ कोशिकाएँ; $\log_{2}10^{9} \approx 30$
द्विगुणन।
\textbf{2.} $30\times 100 = 3000$ दिन, लगभग $8$ वर्ष।
\textbf{3.} \qty{1}{kg} $\approx 10^{12}$ कोशिकाएँ: दस और द्विगुणन,
$1000$ दिन, तीन वर्ष से कम --- पकड़ तक के समय का एक-तिहाई।
\textbf{4.} आरंभिक द्विगुणन-काल लगभग $60$ दिन रहा होगा, या \omterm{def:b3:cancer-biology:hallmarks}{अर्बुद} छोटा
होने पर तेज़ी से बढ़ा और बड़ा होते-होते धीमा पड़ गया (यही सामान्य
प्रतिरूप है, क्योंकि आपूर्ति और जगह चुक जाती हैं)।
\textbf{5.} $2$ दिन का कोई चक्र समष्टि को हर $2$ दिन में दोगुना कर
देता; $100$ का द्विगुणन-काल बताता है कि शुद्ध वृद्धि जन्म-दर का
$2\,\%$ है: जन्मी लगभग हर कोशिका के सामने एक मरती है, या थोड़ी ही
कोशिकाएँ चक्र में हैं --- जन्म और मृत्यु लगभग बराबर, और वृद्धि उनका
छोटा अंतर।
\textbf{6.} औषधें चक्र में चलती कोशिकाएँ मारती हैं, इसलिए जिस \omterm{def:b3:cancer-biology:hallmarks}{अर्बुद}
का थोड़ा-सा अंश चक्र में हो वह प्रति बारी कम अंश खोता है; पर तेज़
\omterm{def:b3:cancer-biology:hallmarks}{अर्बुद}, अधिक खोकर भी, बारियों के बीच फिर बढ़ जाता है। संतुलन में तेज़
\omterm{def:b3:cancer-biology:hallmarks}{अर्बुद} (श्वेतरक्तताएँ, लसीकार्बुद, वृषण-कर्कट) ही वे हैं जिन्हें
रसायन-चिकित्सा ठीक करती है।
\textbf{7.} अनुपात $300$, आयु-अनुपात $80/30 = 2.67$ पर: $k - 1 =
\ln 300/\ln 2.67 = 5.8$, $k \approx 7$।
\textbf{8.} $N(ut)^{k} = 10^{8}\times (7\times 10^{-5})^{7} \approx
10^{-21}$: $1$ प्रति $3$ के आसपास के जीवन-काल जोखिम के सामने बेतुका।
छूट रहा है: हर आघात के बाद जोखिम वाली कोशिकाएँ गुणा करता क्लोनी फैलाव,
अस्थिरता से बढ़ी उत्परिवर्तन-दरें, कहीं अधिक विभाजन (दशकों तक चक्र
लगाती स्तंभ कोशिकाएँ), और सचमुच कम चालक।
\textbf{9.} $g^{k-1} = 100^{6} = 10^{12}$ से गुणा करने पर $10^{-9}$
मिलता है --- सात घटनाओं के साथ अब भी नगण्य; चार घटनाओं और फैलावों के
साथ $10^{8}\times (7\times 10^{-5})^{4}\times 10^{6} \approx 2\times
10^{-3}$, यानी सही कोटि। असली चित्र थोड़े चालकों और बड़े फैलावों का है।
\textbf{10.} $(ut)^{k-1}/(ut)^{k} = 1/(ut) = 1/(10^{-6}\times 40) =
2.5\times 10^{4}$: चालीस पर किसी वाहक की घटना-दर छिटपुट दर से दसियों
हज़ार गुना है --- कुलगत \omterm{def:b3:cancer-biology:hallmarks}{कर्कट} जवानी में चोट करते हैं।
\textbf{11.} एक घटना दस गुना तेज़: घटना-दर $\times 10$; दो:
$\times 100$। प्रेक्षित पंद्रह से पच्चीस गुना जोखिम सुझाता है कि धुआँ
एक से अधिक चरण पर काम करता है।
\textbf{12.} छोड़ने पर बची हुई घटनाओं के लिए $u$ फिर गिर जाता है,
इसलिए घटना-दर उतनी तीखी नहीं चढ़ती; पर कोशिकाओं में पहले से जमा आघात
बने रहते हैं, इसलिए पूर्व-धूम्रपायी की कोशिकाएँ कभी न धूम्रपान करने
वाले की तुलना में \omterm{def:b3:cancer-biology:hallmarks}{कर्कट} के अधिक पास हैं और अतिरिक्त जोखिम बना रहता है,
जमा हुआ, मिटा हुआ नहीं।
\textbf{13.} $L = \sqrt{2\times 2\times 10^{-9}\times 0.05/0.03} =
\qty{82}{\micro m}$।
\textbf{14.} व्यास $2L \approx
\qty{160}{\micro m}$ तक पूरी तरह ऑक्सीजनित। \qty{2}{mm} की गोलिका का
जीवित किनारा \qty{82}{\micro m} है: जीवित अंश $1 - (918/1000)^{3} = 0.23$।
\textbf{15.} बीचोबीच की कोशिकाएँ किसी वाहिनी से \qty{75}{\micro m} दूर
हैं, यानी ठीक $L$ के भीतर --- सिद्धांततः ऑक्सीजनित, पर \omterm{def:b3:cancer-biology:hallmarks}{अर्बुद} की
वाहिनियाँ कम ऑक्सीजन ढोती हैं और रुक-रुककर बहती हैं, इसलिए बीचोबीच की
कोशिकाएँ अल्पऑक्सी हैं। अल्पऑक्सी कोशिकाएँ विकिरण के प्रति दो से तीन
गुना अधिक प्रतिरोधी हैं (क्षति को स्थायी बनाने के लिए ऑक्सीजन चाहिए),
और अल्पऑक्सीयता p53-निर्भर \omterm{def:b3:cell-cycle-apoptosis:apoptosis}{एपॉप्टोसिस} चलाती है, इसलिए वह उन्हीं
कोशिकाओं को चुनती है जिन्होंने \omterm{def:b3:cancer-biology:genes}{p53} खो दिया है।
\textbf{16.} सघनता आधी करने से वाहिनियाँ $\sqrt{2}$ से दूर हो जाती हैं:
\qty{212}{\micro m} की दूरी पर, और बीचोबीच की कोशिकाएँ
\qty{106}{\micro m} दूर, यानी $L$ से परे: वे मर जाती हैं या, अल्पऑक्सी
बचकर, अधिक उग्र और रक्तवाहिनीजनक हो जाती हैं।
\textbf{17.} $30/2 = 15$ गुना अधिक ग्लूकोज़। इसलिए \omterm{def:b3:cancer-biology:hallmarks}{अर्बुद} ग्लूकोज़ के
अनुरूप सांद्र कर लेते हैं, और कोई पॉज़िट्रॉन-उत्सर्जी
फ़्लोरोडिऑक्सीग्लूकोज़ स्कैन उन्हें जगमगा देता है।
\textbf{18.} भूखे \omterm{def:b3:cancer-biology:hallmarks}{अर्बुद} विसरण-सीमित आकार तक सिकुड़कर सुप्त और
अल्पऑक्सी बने रहते हैं, या पहले से मौजूद वाहिनियाँ अपना लेते हैं;
अल्पऑक्सीयता उग्र क्लोन चुनती है। पर \omterm{def:b3:cancer-biology:hallmarks}{अर्बुद} की अस्तव्यस्त
वाहिनी-सज्जा की छँटाई अंतराली दाब घटा देती है और बहाव एकसमान कर देती
है, इसलिए रसायन-चिकित्सा बेहतर पहुँचती है --- संयोजन वहाँ काम करता है
जहाँ कोई अकेला नहीं।
\textbf{19.} $10^{9}\times 10^{-2n} < 1$: $n > 4.5$, पाँच बारियाँ।
\textbf{20.} $21$ दिन $= 0.21$ द्विगुणन, पुनर्वृद्धि $\times 1.16$;
प्रति चक्र शुद्ध गुणक $0.0116$; $\ln 10^{-9}/\ln 0.0116 = 4.65$: फिर भी
पाँच चक्र।
\textbf{21.} $10^{-7}\times 10^{9} = 100$ प्रतिरोधी कोशिकाएँ
प्रत्याशित; $P_{0} = e^{-100} \approx 0$: प्रतिरोध निश्चित है।
\textbf{22.} $\mu_{1}\mu_{2}N = 10^{-14}\times 10^{9} = 10^{-5}$;
$P_{0} = 0.99999$। $10^{12}$ कोशिकाओं पर $10^{-2}$, $P_{0} = 0.99$।
\textbf{23.} $10^{6}$ कोशिकाएँ: एक औषध $\mu N = 0.1$, $P_{0} = 0.90$; दो
औषधें $10^{-8}$, $P_{0} \approx 1$। समष्टि छोटी हो तब उपचार कीजिए,
और संयोजनों से कीजिए।
\textbf{24.} T~कोशिकाएँ एक साथ अनेक नवप्रतिजनों पर वार करती हैं, इसलिए
कोई एक उत्परिवर्तन उनसे बच नहीं निकलता; बचने के लिए स्वयं
प्रतिजन-प्रस्तुति खोनी पड़ती है (MHC, $\beta_{2}$-सूक्ष्मग्लोब्युलिन)
या T~कोशिकाएँ दबानी पड़ती हैं, जो एक भिन्न और दुर्लभ वर्ग की घटना है।
वही ऊँचा उत्परिवर्तन-भार, जो किसी \omterm{def:b3:cancer-biology:hallmarks}{अर्बुद} को किसी काइनेज़ निरोधक का
प्रतिरोध करना निश्चित बना देता है, उसे अनेक नवप्रतिजन देता है:
असंगति-मरम्मत-अपर्याप्त, पराबैंगनी- और धुआँ-प्रेरित \omterm{def:b3:cancer-biology:hallmarks}{अर्बुद} सबसे अच्छी
अनुक्रिया देते हैं।
\textbf{25.} पकड़ तक लगभग $8$ वर्ष; किसी वाहिनी के चारों ओर
\qty{82}{\micro m} गहरा जीवित \omterm{def:b3:cancer-biology:hallmarks}{अर्बुद}; $P_{0} = 0.99999$ कि
\qty{1}{cm^3} के किसी \omterm{def:b3:cancer-biology:hallmarks}{अर्बुद} में दो-औषध संयोजन को कोई प्रतिरोधी कोशिका
न मिले।
\end{solution}
